\subsection{Boundary bundles and mapping-cylinder parametrizations}
\label{subsec:boundary-bundle-rulings}
All spaces and maps in this subsection are definable in the fixed
o-minimal expansion of a real closed field $K$, and the objects are
closed and bounded.  A bordant pair $(E,B)$ has a collared boundary
subset $B$.  A boundary-bundle object consists of $(E,B)$, a core
space $Y$, and a specified locally trivial definable bundle
$\varphi:B\to Y$.  Its morphisms are homeomorphisms of the pairs
and of the cores whose boundary restrictions satisfy the bundle
square.  The unrigidified collapse identifies each boundary fibre
with its point in $Y$ and remembers only the resulting pair $(X,Y)$.

Collapsing a fibre can lose the limiting directions needed to lift a
homeomorphism back to the boundary. Specifying the product
parametrization makes the extension of a boundary-compatible map
unique; cutting off the open inner tube then gives the inverse
construction.

\begin{proposition}[Boundary collapse and product parametrizations]
\label{prop:boundary-bundle-retained-equivalence}
The unrigidified collapse functor is not full as an ordinary functor,
already for a one-point core.  In the stipulated definable simplicial
categories it also fails to be a Dwyer--Kan equivalence.

For the corrected target, an object consists of closed bounded definable sets
$(X,Y,V,B)$ with $Y\subset V\subset X$ and
$B\subset V\setminus Y$, a specified locally trivial definable bundle
$\varphi:B\to Y$, and a specified homeomorphism
$\eta:M_\varphi\to V$ that is the identity on the core $Y$ and
identifies $B\times\{1\}$ with the outer face $B$.
Require that $V\setminus B$ be open in $X$ and that $B$ have an
exterior collar in $X\setminus(V\setminus B)$.
A target morphism preserves these marked sets and the full
presentation: on the tube it has the form
\[
 [b,t]\longmapsto[h_B(b),t],
 \qquad h_Y\varphi=\varphi'h_B,
 \qquad 0<t\leq1.
\]
Here $h_B$ is its restriction to the outer face and $h_Y$ its
restriction to the core.
With these object data and morphisms, attachment
\[
 (E,B,Y,\varphi)\longmapsto E\cup_B M_\varphi
\]
is an equivalence and induces isomorphisms of the definable
simplicial mapping spaces.
\end{proposition}
\begin{proof}
For ordinary fullness take $Q=[-1,1]^2$, $B=\partial Q$, and
the cylinder $B\times[0,1]$ with its inner boundary collapsed to
one point.  The quotient is $Q$, with quotient map
\[
 q(b,t)=tb.
\]
The semialgebraic homeomorphism
$h:Q\to Q$, $h(x,y)=(x,y^3)$, fixes the distinguished point.
Any lift to a cylinder homeomorphism is forced away from the
inner boundary, because $q$ is a homeomorphism there.
On the right edge, for $b=(1,s)$ and $t>0$, the forced lift is
\[
 ((1,s),t)\longmapsto((1,t^2s^3),t).
\]
Indeed the maximum norm of $(t,t^3s^3)$ is $t$.
Continuity at $t=0$ would send every point of that inner edge to
the same point $((1,0),0)$, which excludes a homeomorphism.
The argument uses polynomial identities and order inequalities,
so it works over every $K$ as stated.

The square homeomorphism is itself joined to the identity by
\[
 h_u(x,y)=\bigl(x,(1-u)y+uy^3\bigr),\qquad0\le u\le1.
\]
Each map in this definable family is a homeomorphism: its second
coordinate is strictly increasing from $-1$ to $1$.
For the enriched conclusion we use the six-torus obstruction of
Theorem~\ref{result-5}.  Here the boundary bundle over a point is
automatically locally trivial, and equisingularity along that core
places no additional restriction.  To recall the precise mapping
space comparison, put $T=(S^1)^6$ and
\[
 \begin{split}
 G&=\operatorname{Homeo}_{\rm def}
       (T\times I;T\times\{0\}\text{ setwise}),\\
 H&=\operatorname{Homeo}_{\rm def}(CT;o).
 \end{split}
\]
Restriction to the outer boundary maps both groups to
$\operatorname{Homeo}_{\rm def}(T)$ and has sections given by
constant extension and by coning.  These sections identify the
underlying simplicial spaces with the product of the restriction
kernel and $\operatorname{Homeo}_{\rm def}(T)$.  The kernel for
$G$ is the group fixing the upper boundary, which becomes
$P_{\rm def}(T)$ after interval reversal.  The kernel for $H$
fixes the base and vertex and contracts by
Theorem~\ref{result-1}.  Boundary collapse commutes with these
restriction maps.  Its homotopy fibre is therefore
$P_{\rm def}(T)$, whose component group is nonzero by
Theorem~\ref{result-3}.  Thus the induced mapping-space map is
not a weak equivalence.  This is the required Dwyer--Kan obstruction
inside the one-point-core subcategory.

For the positive assertion attach the full mapping cylinder along
its outer face.  The product collar on the exterior side comes
from the bordant pair; the mapping-cylinder side supplies the
inner product and its limiting core.  Thus the specified tube is
strict.  Local triviality of $\varphi$ identifies its neighborhood
of the core, locally over $Y$, with the product of an open subset
of $Y$ and the cone on the bundle fibre.

Conversely, from a target object $(X,Y,V,B)$ with the stipulated
strict mapping-cylinder description recover the exterior
\[
 E=X\setminus(V\setminus B).
\]
The complement is closed, because $V\setminus B$ is the open
inner part of the strict tube.  Its outer collar makes $(E,B)$
a bordant pair.  Reattaching $M_\varphi$ reconstructs the original
target with its given markings and parametrization.

A resolved morphism $(h_E,h_Y)$ has boundary restriction $h_B$.
It extends to the tube by the displayed product formula, and to
the core by $h_Y$.  At $t=0$ the compatibility equation
$h_Y\varphi=\varphi'h_B$ makes this well-defined and continuous
on the quotient.  At $t=1$ it agrees with $h_E$ on the glued
outer face.  Applying the same construction to the inverse data
proves that the extension is a homeomorphism.  A target morphism
with the stipulated product formula restricts to exactly those
resolved maps, so restriction and extension are inverse on
morphism sets and commute with composition.

The same formulas apply to a joint definable family over every
standard simplex $\Delta^p_K$.  Continuity at the core can also
be checked in the closed bounded model consisting of the points
$(tb,\varphi(b),t)$ together with the entire core copy
$\{(0,y,0):y\in Y\}$. The output first coordinates are $t h_B(b)$,
with $h_B$ uniformly bounded, and the core coordinate is
$h_Y(\varphi(b))$, with the parameter included in both maps.
These formulas are jointly continuous at $t=0$.  Gluing along
the closed outer face gives the required family, and the same
argument applies to the inverse.  Parameter pullback leaves
the formulas unchanged.
Thus extension and restriction commute with faces and
degeneracies and are inverse on the simplicial mapping spaces.
If $B$ is empty, the construction is the disjoint union $E\amalg Y$,
and the same restriction and extension statements hold.

The full product condition is necessary for this particular
inverse construction.  For instance, with $B=[-1,1]$ and
one-point core, put $a(t)=t(1-t)$.  The map
\[
 [b,t]\longmapsto
 [(1-a(t))b+a(t)b^3,t]
\]
is a definable radius-preserving tube automorphism, equal to the
identity on its outer face and core, and nonidentity inside.
Strict monotonicity in $b$ and the cone topology give its
continuous inverse.  It extends by the identity on the exterior.
Consequently retaining only the radius would leave distinct
target morphisms with the same resolved restriction.
\end{proof}
